\section{The proof of a goal}\label{sec:qed}

\Cref{sec:prooflog} turned Sylvan's log into \emph{facts} about decision
diagrams. This section shows how the tactic turns those facts into a
proof of the goal itself, so that a file ends with \code{Qed}:
\begin{verbatim}
Example s1 : strong_sim step step (par A B) (par B A).
Proof. sim. Qed.
\end{verbatim}
The tactic decides the goal, Sylvan logs the decision (the
\emph{certificate}), and the tactic builds from it a proof that Rocq's
kernel checks at \code{Qed}. The proof uses no axiom: \code{Print
Assumptions s1} says ``Closed under the global context''
(\file{test/rocq/LogTest.v}). The section first defines what it needs
from Rocq (\cref{sec:qed:rocq}), then the example
(\cref{sec:qed:example}), the theorem the proof applies
(\cref{sec:qed:theorem}) and how each of its hypotheses is proved
(\cref{sec:qed:parts}); then how the pieces are put together
(\cref{sec:qed:assembly}), what Rocq computes (\cref{sec:qed:computes}),
how large goals were made to fit (\cref{sec:qed:fit}), the cost
(\cref{sec:qed:cost}) and how to try it (\cref{sec:qed:try}).

\subsection{What this section needs from earlier sections}\label{sec:qed:terms}

Each term below is defined, with an example, where the cross-reference
points; here is what each means for this section.
\begin{itemize}
\item A \emph{diagram} is a set of vectors of numbers, stored as a tree of
  \emph{nodes}: a node holds a value, a child \emph{down} (the rest of the
  vectors with that value) and a child \emph{right} (the vectors with a
  larger value there); the nodes reached by following right children
  from a node are its \emph{siblings} (\cref{sec:prooflog:background}).
\item A \emph{relation} is a diagram of \emph{tuples}: per position of a
  vector, a code says whether the step copies the value, reads it, or
  reads it and writes a new one, and the tuple holds the values read and
  written (\cref{sec:prooflog:images}).
\item The \emph{mover} is the system whose moves must be answered, the
  \emph{answerer} the one that answers them
  (\cref{sec:prooflog:background}).
\item A \emph{summand} is one way a state of a given shape can step: a
  chain of rules from the shape to a step of one of its slots
  (\cref{sec:containers}).
\item A \emph{fact} is one lemma of \file{theories/DdTree.v} applied to
  nodes and to other facts; it states what one of Sylvan's computations
  produced (a union, a difference, an image), and the kernel accepts it
  by typechecking (\cref{sec:prooflog:facts}). A \emph{chunk} is a batch
  of facts with the nodes they name, translated from the log together.
\end{itemize}

\subsection{What this section needs from Rocq}\label{sec:qed:rocq}

\paragraph{Terms, types, the kernel.} In Rocq a proof is a \emph{term}
whose \emph{type} is the statement it proves: a function taking a proof
of $A$ to a proof of $B$ proves $A \to B$. The \emph{kernel} is the small
part of Rocq that checks a term against its type (it \emph{typechecks}
it); nothing else is trusted. A \emph{tactic} is a program that builds
a term while a proof is in progress; \code{Qed} hands the finished term
to the kernel. The plugin's tactic builds the whole term at once and
gives it to Rocq without checking it (\code{exact_no_check}): the
kernel's check at \code{Qed} is the only one that counts.

\paragraph{Conversion.} While typechecking, the kernel may need two types
to be equal that are written differently, and it accepts them when one
\emph{computes} to the other: unfolding a definition, applying a function
to its argument, taking the matching branch of a \code{match}. This is
\emph{conversion}. \emph{Example:} if \code{val} looks a number up in a
table of values, a proof about \code{val 3} serves for the value stored
at 3, because \code{val 3} computes to it. Conversion is the only
computation the plugin's proofs rely on (\cref{sec:qed:computes}).

\paragraph{Definitions and \code{let}.} A \emph{constant} is a named
global definition: a body and its type. It is \emph{transparent} when
conversion may unfold it (a datum, such as a table of values, must be:
its contents matter), and \emph{opaque} when it may not (a proof: only its
statement matters, so nothing ever needs to look inside it). Inside a
term, \code{let x := b : T in c} names \code{b} as \code{x} within
\code{c}; the kernel checks \code{b} once, however many times \code{c}
uses \code{x}.

\paragraph{Constants made by a tactic.} Rocq's tactic \code{abstract}
proves a sub-goal and declares the proof as an opaque constant of its own,
a \emph{side effect} of the proof in progress: the kernel checks it at
once, the proof's term only names it, and at \code{Qed} it is kept as a
global constant (\code{s1_subproof}). The plugin declares its pieces the
same way, named after the lemma (\code{s1_mb_case3} and so on).

\paragraph{Elaboration.} Rocq's \emph{elaborator} turns text written in
Gallina (Rocq's language) into a term, inferring what the text leaves
out: the types of the arguments marked \code{_}, the type of a
\code{match}'s result. The plugin writes the parts of a proof that are
specific to a goal as such text and has Rocq elaborate it.

\subsection{The running example}\label{sec:qed:example}

\file{test/rocq/LogTest.v} defines processes \code{pr} (\code{nil},
\code{pre x k}: the action \code{x} then \code{k}, and \code{par p q})
and their steps \code{step}:
\begin{itemize}
\item \code{s_pre}: \code{pre x k} steps to \code{k} with label \code{x};
\item \code{s_parl}: if \code{p} steps to \code{p'}, then \code{par p q}
  steps to \code{par p' q};
\item \code{s_parr}: the same on the right.
\end{itemize}
With \code{A = pre (Some a) (pre None (pre (Some a) nil))} and
\code{B = pre (Some b) nil}, the goal says that \code{par B A} can match
every step of \code{par A B} (\code{strong_sim}). The engine numbers the
values it meets (1 is \code{A}, 2 is \code{B}, 3 is \code{pre None
(pre (Some a) nil)}, 4 is \code{nil}, 5 is \code{pre (Some a) nil}) and
the shapes (shape 1 is \code{par} with two slots). A state is then a
\emph{vector}: \code{par A B} is $[1; 1; 2; 0; \dots]$ (shape 1, slot 0
holds value 1, slot 1 holds value 2, unused positions 0). A pair of
states is one vector, the answering system's positions first
(\cref{sec:prooflog:fit}).

\subsection{The theorem}\label{sec:qed:theorem}

\file{theories/DdGame.v} proves once, for any systems, this theorem
(\code{cert_strong_sim}; the weak and bisimilarity versions are alike):
if
\begin{enumerate}
\item the \emph{game on vectors} holds: every move from a pair of the
  diagram $z$ by a relation of the mover is answered into $z$ by a
  relation of the answerer (\code{VSSim}; \code{VWSim} for weak);
\item each system's states are \emph{encoded} as vectors (\code{EL},
  \code{ER}), and the left system's vectors all have its width
  (\code{LenL});
\item the mover's relations are \emph{complete}: every step of an encoded
  state of $z$ is a step of one of its relations (\code{CompleteR});
\item the answerer's relations are \emph{sound}: every step of one of its
  relations is a step of its states (\code{SoundL}),
\end{enumerate}
then any two states whose vectors form a pair of $z$ are related. The
tactic proves the four hypotheses for its goal, and that the start pair
is in $z$. Here $z$ is the diagram of pairs the engine's game ended with
(\cref{sec:prooflog:background}); a \emph{relation} is a diagram of
tuples, each a step of one rule with one label
(\cref{sec:prooflog:images}).

\subsection{The four hypotheses}\label{sec:qed:parts}

\paragraph{1. The game} comes from the log's facts (\cref{sec:prooflog}):
\file{src/game_proof.ml} applies the lemmas of \file{theories/DdGame.v}
to the facts in the order the logged check ran. A sub-proof used twice
(an image's property, a round of the silent closure) is a \code{let}:
written out each time, the closure's rounds make the term grow
exponentially.

\paragraph{2. The encoding} (\code{DdEnc.enc}). Each shape has a
predicate: for shape 1,
\begin{verbatim}
fun x rest => exists u0 u1, rest = [u0; u1; 0; ...]
                         /\ x = par (val u0) (val u1)
\end{verbatim}
where \code{val} looks a value's number up in a binary tree of the values
(\code{DdEnc.val}; the kernel follows one branch per binary digit). A
vector \code{s :: rest} encodes \code{x} when the predicate of shape
\code{s} holds of \code{x} and \code{rest}; the predicates are kept in a
tree too, so the first number picks the shape. \code{par A B} and
$[1; 1; 2; 0; \dots]$: take $u_0 = 1$, $u_1 = 2$; both equations hold by
conversion (computing \code{val}).

\paragraph{3. Soundness.} The answerer's relations are logged backwards
(\cref{sec:prooflog:images}); each tuple, read forwards, must be a real
step. For the relation of \code{s_parr} with label \code{Some a} of the
system \code{par B A}, a tuple says ``slot 1 holds 1 and becomes 3''.
One lemma per relation and pair of shapes says: for any numbers $c$, $w$,
\emph{if} \code{par (val u0) (val c)} steps with \code{Some a} to
\code{par (val u0) (val w)} for every $u_0$, \emph{then} the tuple
$(c, w)$ is sound. Its proof takes the vector apart (the shape, then the
template) and rewrites the slot it reads. Each tuple then applies the
lemma to its numbers and to a \emph{derivation}, the rules applied:
\begin{verbatim}
fun u0 => s_parr (val u0) A (Some a) (pre None (pre (Some a) nil))
                 (s_pre (Some a) (pre None (pre (Some a) nil)))
\end{verbatim}
The engine's rule interpreter knows the derivation
(\code{Learn.derive_chain}): it is the same interpretation that found the
step. \code{AllT} (\file{theories/DdEnc.v}) walks the relation's diagram,
one tuple per leaf: \code{AllT d Q} says that \code{Q} holds of every
tuple of the diagram \code{d}, and is proved node by node (\code{allt_0},
\code{allt_1}, \code{allt_nd}).

\paragraph{4. Completeness} is the hard direction: every step of the
mover's state is a tuple. Take a state of shape 1,
\code{par (val u0) (val u1)}, and a step \code{T} from it.
\begin{itemize}
\item \emph{The case lemma}, one per LTS, says which rule made a step:
  it is a match on the step (no tactic), one disjunct per rule, each with
  the rule's variables, its premises, and equations ``the source is the
  rule's source'', ``the label is the rule's label'', ``the target is
  the rule's target''.
\item \emph{Small inversion} \citep{monin-shi13} removes the equations. For \code{s_pre}, the
  equation \code{par (val u0) (val u1) = pre x k} is impossible: a match
  on it, with \code{par} giving \code{True} and anything else
  \code{False}, proves \code{False} from \code{I}. For \code{s_parl},
  the equation \code{par (val u0) (val u1) = par p q} replaces \code{p}
  and \code{q} by the slots' values in every remaining hypothesis.
\item \emph{Coverage.} The premise is now a step of \code{val u0}, a
  value we know nothing about. The engine learned the steps of slot 0
  for every value slot 0 takes in $z$ (its \emph{projections}: 1, 3, 4
  and 5), and logged that $z$ has no other: the image of $z$ by the
  shape's filter (complete) lies within its image by the projections
  (sound), \code{DdEnc.cov_of}. So the proof splits on $u_0 \in \{1, 3,
  4, 5\}$, one case (a \emph{leaf}) per projection, \code{AllT} over the
  projections' diagram. When the remaining hypotheses inspect only some
  of the slots (in CADP, a process's program point, not its variables),
  the projections that agree on those slots share one lemma over the
  other slots' numbers, and each projection is that lemma applied to its
  numbers and to its steps' \emph{witnesses}, terms built without
  elaboration.
\item \emph{Step lemmas.} In the case $u_0 = 1$, the premise is a step of
  the closed term \code{A}. Its step lemma, proved the same way, says
  that \code{A}'s only step is \code{Some a} to \code{pre None (pre (Some
  a) nil)}. A step lemma is per \emph{skeleton}: positions no rule
  inspects are holes. The rule \code{s_pre} inspects neither argument of
  \code{pre}, so one lemma serves every \code{pre x k}: its only step is
  \code{x} to \code{k}. In CADP the processes' state is never inspected
  by the rules' patterns, so one lemma per program point serves every
  state and resource.
\item \emph{The witness.} The proof's goal is then a step of a known
  label and target, and the relation's tuple $(1, 3)$ for label
  \code{Some a} matches it. The witness says so: the tuple is in the
  relation's diagram (its \emph{membership}, \code{In_t r t}), what it
  reads agrees with the state (\code{Reads}, each equation by
  \code{eq_refl} on numbers), and the target is encoded by the vector
  after the step (\code{post}). Memberships come from the log
  (\cref{sec:qed:packs}).
\end{itemize}
Some cases are impossible and are refuted from the encoding: a
projection of another shape (by the shape's number); a rule needing a
slot to start with a constructor no value of its type starts with (in
Proc, a reassociation needs \code{tpar} inside a slot;
\code{DdEnc.val_all} proves no value has it); a shape no summand has
($z$ holds none: an empty image, logged).

\subsection{How the proof is put together}\label{sec:qed:assembly}

In the order the tactic works (\file{src/user_tactics.ml},
\code{prove}; \file{src/lts_proof.ml}, \code{qed_term}):
\begin{enumerate}
\item \textbf{Decide and log.} The engine decides the goal and repeats
  the game's last check with Sylvan's log on (\cref{sec:prooflog:log}).
  The log stays in Sylvan's memory: no file is written.
\item \textbf{The plan.} The engine records what it logged and what the
  proof needs from it: the shapes and values, each relation's label,
  codes, diagram and rules, each summand's projections
  (\code{Learn.Plan}).
\item \textbf{The pieces.} The tactic builds the encodings, the case
  and step lemmas, the coverages and the soundness of each relation
  (\cref{sec:qed:parts}). Each piece is a constant, declared as
  \code{abstract} declares a sub-proof (\cref{sec:qed:rocq}): a proof
  opaque, a datum transparent, or written in place when it is a dozen
  nodes or fewer. When a later piece's text is elaborated, a proof it uses
  is visible as its statement only (it is never unfolded), a datum with
  its body. A piece that mentions a fact of the certificate (a relation's
  diagram, the list of relations) cannot know the fact yet, since the log
  is read last: it takes the fact as an argument, and is used applied to
  it.
\item \textbf{The facts.} The tactic asks the log's translator for what
  the pieces and the game need: an image's property, a union, a
  difference, a node, a tuple's membership (each a \emph{request}). The
  translator answers in \emph{chunks}, each a list of nodes and facts
  with nodes of its own (\cref{sec:prooflog:facts}): the main chunk, then
  one chunk per relation diagram the coverages name, then one chunk per
  coverage's memberships (\cref{sec:qed:packs}). The lines go from the
  translator to the chunk's reader as they are written: no file.
\item \textbf{The memberships} of each coverage become one constant
  (\cref{sec:qed:packs}), checked by the kernel at once.
\item \textbf{The term.} The main chunk's nodes and facts become
  \code{let}s; inside them, the game proved from the facts, the coverage
  of each summand (its logged facts), and the theorem applied to the
  pieces, each piece named by its constant applied to its facts. The
  tactic closes the goal with that term; the kernel checks it at
  \code{Qed}, with the constants it names already checked.
\end{enumerate}

\subsubsection{Memberships, packed}\label{sec:qed:packs}

A coverage of three-process CADP has 21,903 leaves, each with a witness
whose tuple must be in its relation: thousands of distinct memberships
\code{In_t r t}, one fact each in the log. Two things make that costly.
\begin{itemize}
\item \emph{An application's cost.} A constant over $k$ facts is used
  applied to $k$ arguments, and the kernel types such an application in
  time proportional to $k^2$: for each argument it rebuilds the
  argument's expected type under the substitution of all the arguments
  before it (\code{Typeops.type_of_apply}). With $k$ in the thousands, a
  single application took minutes.
\item \emph{The memberships' size.} Proved one by one, every membership
  walks its relation's diagram from the top
  (\cref{sec:qed:fit:members}).
\end{itemize}
So a coverage's memberships are \emph{packed} into one fact. Their
statements are numbered $1, \dots, k$ and kept in a binary tree by
number, a \emph{trie} (\code{DdEnc.trie}): a datum, \code{MT}. One fact,
\code{TAll MT}, says that every statement in the trie holds. A leaf takes
its membership by its number $i$ with the lemma \code{DdEnc.tall_get}:
from \code{TAll MT}, the entry of \code{MT} at $i$. The kernel finds that
entry by conversion, following one branch per binary digit of $i$, as it
finds a value by \code{val}. The coverage then takes one argument for
all its memberships.

\emph{Example.} Three memberships, numbered 1 to 3:
\code{In_t r [1;3]}, \code{In_t r [2;3]}, \code{In_t r [5;4]}. The trie
holds the first at its root (1), the second down the left (binary 10),
the third down the right (11). The leaf that needs the third uses
\code{tall_get MT 3 H}, of type ``the entry at 3'', which converts to
\code{In_t r [5;4]}.

The pack itself is proved apart from the main term, so that the kernel
checks it while the proof is built, not inside the term it checks at
\code{Qed}. The translator writes a chunk for each relation diagram the
packs mention: its nodes only; the diagram becomes a transparent
constant, \code{RR}, its nodes as \code{let}s. Then a chunk for each
pack: the membership facts and the relations' nodes, its own. The pack
becomes an opaque constant: the chunk's nodes and facts as \code{let}s,
then the proof of \code{TAll MT} over the chunk's own diagrams, given the
statement \code{TAll MT} over the \code{RR} constants. The two statements
differ only in which copy of each diagram they name, and the kernel
compares them once, diagram against diagram, not entry by entry. In the
main term, the coverage's pack argument is the pack's constant.

\subsection{What Rocq computes}\label{sec:qed:computes}

Only conversions, while typing the term: a value's number to its term
(\code{val}); a coverage's membership by its number in the trie of the
coverage's memberships (\code{tall_get}, the same lookup; accepted on 11
November 2026 as of the same kind as \code{val}); with \code{MeBi
Prove}, a node by its number in its height's trie (\code{val}; accepted
on 12 November 2026) and a fact of an earlier segment by its key
(\code{tall_get}) (\cref{sec:qed:prove}); a vector after a
tuple, and a projection of it, on lists of known length (\code{post},
\code{Reads}, \code{projv}); the user's own functions in a rule (CADP's
\code{do_act}, CCS's \code{def}); a decidable side condition's boolean
(CCS's \code{allowed}), through its \code{Decidable} instance; and a
diagram of the log against the same diagram built from its tuples (a
coverage's projections, a relation's tuples), compared node by node
where the two meet. No boolean about the certificate is proved by
\code{eq_refl}: properties of code lists are inductive predicates proved
by their constructors (\code{NoWO}, \code{Within}, \code{Shifted},
\code{RoOnly}).

\subsection{Making large goals fit}\label{sec:qed:fit}

Each measure below answers a cost that was measured: with timestamps
the plugin prints (\code{MEBI_CERT_STATS=1}), and once with a stack
sample of the running process under \code{gdb}. Two-process CADP has
12,027 pairs, three-process CADP 7.9 million.

\paragraph{The pieces as constants.} With the pieces as \code{let}s of
one term, the kernel checks everything at \code{Qed}, and the elaborator
can unfold a proof while it elaborates the texts that use it. As opaque
constants, each piece is checked once when declared, and never unfolded.
One-process CADP: \code{Qed} 0.11\,s instead of 0.29\,s.

\paragraph{Subterms written once.} A term is a tree in which a subterm
written twice is checked twice. Three places repeated subterms:
\begin{itemize}
\item \code{allt_nd v d r Q pd pr}, the step of \code{AllT} at a node,
  names the node's two subdiagrams \code{d} and \code{r} in full. Along a
  list of 1,000 siblings, the last sibling is named a thousand times. Each
  subdiagram is now a \code{let}, named once. A \code{let} has its own
  cost: the kernel types \code{let x := b in c} by substituting \code{b}
  into the type of \code{c}, and if that type names \code{x}, every
  \code{let} enlarges it (448\,s for one coverage of three processes). So
  the proof is given its statement (\code{AllT D Q}, with \code{D} the
  diagram's constant) inside the \code{let}s, by an identity function of
  that type, and the substitutions stay small.
\item Lists of numbers (the rest of a vector, at each level of a
  membership) share their tails, and each variable bound by a \code{let}
  is one term, however many times it is used (two processes: the facts
  18\,MB down to 4\,MB).
\item Substituting a large term under binders copies it, and so does
  shifting it to be used under more binders: a pack, substituted where
  the coverage takes it, was copied whole. A pack is now a constant,
  named where it is used.
\end{itemize}

\paragraph{Long applications.} As \cref{sec:qed:packs} explains, the
kernel types an application of $k$ arguments in time proportional to
$k^2$. Before packs, the proof of the shapes (which applies each coverage
to its facts) took 2.5\,s of the kernel's time on two processes, and
did not finish in 20 minutes on three; with packs, it takes under
0.01\,s on two processes.

\paragraph{Memberships, one fact per level.}\label{sec:qed:fit:members}
A membership \code{In_t x w} was proved by following \code{x}: a step
right (\code{in_r}) for each sibling whose value is not $w$'s next, a step
down (\code{in_d}) for the sibling that is. \emph{Example:} in the
diagram of $\{(1,5), (2,5), (3,5)\}$, the vector $(3,5)$ takes two steps
right (past 1 and 2), one down, one down. Each step's fact names the rest
of the vector, so two vectors scanning the same siblings share nothing.
Three processes: 9.4M of the proof's 11.1M facts were such steps right.
Now the scan is its own fact, \code{Sib x y} (\code{y} is \code{x} or a
sibling to its right: \code{sib_refl}, \code{sib_step},
\file{theories/DdTree.v}), which names no vector: every vector scanning
from \code{x} shares it, each step made once. A level is then one fact,
\code{in_sib_d x v d r w} (from \code{Sib x (Nd v d r)} and \code{In_t d
w}, \code{In_t x (v :: w)}). And the steps from a node depend on the node
and the rest of the vector only, so vectors meeting at the same node
(Sylvan shares equal subdiagrams) share the rest. Three processes: 11.1M
facts down to 1.1M.

\paragraph{The engine's own lookups.} Building a leaf looked a tuple's
membership up in a list of all memberships so far: 107\,s for one
coverage of three processes, now a hash table (under a second). The
labels' normal forms are remembered; a target is compared with a tuple's
by conversion, which stops at the first difference, instead of computing
both normal forms.

\paragraph{Memory at \code{Qed}.} The kernel's check of the main term
needs memory in proportion to its facts, about 6\,KB each: the types it
computes, the \code{let}s' types, and \emph{hash-consing} (the kernel
rebuilds a term so that equal subterms are stored once, with a table of
them). Before \code{Qed}, the tactic gives
back what it no longer needs: the engine's sets, Sylvan's tables (grown
to their maximum for the log), and the garbage of building the proof.
OCaml's garbage collector lets the heap grow to a multiple of the live
data (Rocq sets that margin, \code{space_overhead}, to 200\%); the tactic
lowers it to 80\% while it builds a proof, and after a decision of more
than a million pairs to 60\% until the next tactic, so that \code{Qed}
and writing the compiled file need less.

\subsection{Large goals: the command \texorpdfstring{\code{MeBi Prove}}{MeBi Prove}}\label{sec:qed:prove}

\paragraph{Why.} The tactic's proof is one term: the main chunk's nodes
and facts as \code{let}s, and inside them the rest
(\cref{sec:qed:assembly}). The kernel checks it at \code{Qed} in one go,
and its memory grows with the facts it has checked so far: three-process
CADP, 829K facts, from 1.6\,GB to 6.6\,GB. Nothing in that check needs
the whole term at once: a fact needs only the nodes and facts it names.
So the proof can be checked in pieces, each declared as a global
constant: the kernel checks a piece when it is declared, and from then on
Rocq keeps the piece, stored once, but not the work of checking it.

\paragraph{Use.} The command takes a name and a statement, one of the
four relations between closed terms, as the tactics' goals are:
\begin{verbatim}
MeBi Prove wsim_lts_3 : weak_sim lts spec_lts c3 (make_spec 2).
\end{verbatim}
It decides the statement and proves it from its certificate as the
tactics do, then declares the lemma \code{wsim_lts_3}, used like any
other (\file{examples/CADP/ThreeProcesses.v}; \file{test/rocq/LogTest.v}
prints its assumptions: none). A negative verdict is an error, and no
lemma is declared.

\paragraph{The nodes, by height.} A diagram's leaves, \code{D0} (the
empty set) and \code{D1}, have height 0; a node \code{Nd v d r} (its
value \code{v}, the diagram \code{d} below it, its right sibling
\code{r}) has height one more than the greater of \code{d}'s and
\code{r}'s. The nodes of each height $h$ are numbered $1, 2, \dots$ and
kept in a trie by number, as a pack's memberships are
(\cref{sec:qed:packs}): a transparent constant, \code{mb_dn}$h$. A node
is written \code{val dd D0 mb_dn}$h$\code{ p}, its number $p$ in its
height's trie: the kernel finds it by conversion, as it finds a value
(\code{val}, \cref{sec:qed:rocq}; this lookup was accepted on 12 November
2026 as of the same kind). A node's children are lower, so they are in
tries declared before its own.

\emph{Example.} The diagram of $\{(1,5), (2,5), (3,5)\}$ has one node
for the second position, $X = $ \code{Nd 5 D1 D0} (height 1), and three
siblings for the first, each with $X$ below it: \code{Nd 3 X D0} (height
2), \code{Nd 2 X (Nd 3 X D0)} (height 3), \code{Nd 1 X (Nd 2 \dots)}
(height 4). So \code{mb_dn1} holds $X$ at 1, \code{mb_dn2} holds
\code{Nd 3 (val dd D0 mb_dn1 1) D0} at 1, and so on: four tries, each a
single entry here. Three-process CADP has 360 heights.

Naming a node by number matters for the facts: a fact's statement names
nodes, and a statement written outside the term that defines its nodes
(the \code{let}s of \cref{sec:qed:assembly}) must name them by something
global. A constant per node would do, but make hundreds of thousands of
global declarations; a trie per height makes one per height.

\paragraph{The facts, in segments.} The translator writes the facts in
an order where each names only facts before it, numbered $0, 1, \dots$.
They are cut into \emph{segments} of 20,000 (\code{MEBI_CHUNK_FACTS}):
facts $0$ to $19{,}999$, then from $20{,}000$, and so on (three
processes: 42 segments). A fact that a later segment or the rest of the
proof names is \emph{exported} by its segment $s$:
\begin{itemize}
\item its statement goes in the segment's trie of statements,
  \code{mb_ds}$s$ (transparent), keyed $1, 2, \dots$;
\item the segment's proof, \code{mb_dp}$s$\code{ : TAll mb_ds}$s$
  (opaque), proves every statement in it: inside it, the segment's facts
  used twice or more are \code{let}s, as in the tactic's term, and the
  proof is given its statement inside them by an identity function, as
  in \cref{sec:qed:fit};
\item a later segment, or the lemma, names the fact as
  \code{tall_get mb_ds}$s$\code{ }$k$\code{ mb_dp}$s$, whose type is the
  entry at $k$, which converts to the statement (\cref{sec:qed:packs}).
\end{itemize}
The rest of the proof (encodings, case and step lemmas, coverages,
packs, soundness) is made as by the tactic, each piece a global constant
instead of a side effect. The lemma's own term names only the pieces:
three processes, checked in 0.4\,s.

\paragraph{Cost.} Three-process CADP: the proof built, every piece
checked as it is declared, in 231\,s (of which the kernel's time: packs
73\,s, segments 52\,s, coverages 37\,s, node tries 19\,s); the lemma
0.4\,s; the file in 279\,s, 2.6\,GB at most. Two processes: the file in
5.8\,s, 0.72\,GB (by the tactic: 6\,s, 0.83\,GB). The time is about the
tactic's (293\,s); the memory is less than half.

\paragraph{Why a tactic cannot do the same.} Each way was tried, and
the reason read in Rocq 9.3's source:
\begin{itemize}
\item \emph{Global constants from a tactic.} \code{rocq compile} runs
  the sentences of a proof on a branch of their own (the branches of its
  document manager, the STM), and \code{Qed} takes the proof's term at
  the end of that branch, then declares the lemma in the state from
  before the proof (\file{stm/stm.ml}, the synchronous \code{Qed}). What
  a tactic declares globally is on the proof's branch: at \code{Qed} the
  term names constants that are gone (``Constant \dots does not appear
  in the environment'').
\item \emph{Constants as side effects} (as \code{abstract} makes,
  \cref{sec:qed:rocq}), the only constants a proof passes on. The kernel
  keeps each with its typed form (a copy of the term carrying what the
  kernel computed while typing it, \code{HConstr}, about three times the
  term: \code{seff_body} in \file{kernel/safe_typing.ml}) until
  \code{Qed} exports them; a global constant's is dropped once it is
  checked. Three processes in pieces: 1.65\,GB to 4.1\,GB while
  declaring them, beyond 7\,GB at \code{Qed}: more than the single term.
\item \emph{A command inside the proof} (written between two tactics)
  declaring the pieces. The STM runs such a command twice: in the proof,
  and at \code{Qed} on the main branch, outside the proof and without
  its goal (only \code{Definition x := body} is run once, its
  environment copied into the proof). So the work is done twice.
\item \emph{As Program's obligations are recorded.} They are kept in the
  state that commands receive (\code{Declare.OblState},
  \file{vernac/vernactypes.ml}); a tactic receives only its goals.
\end{itemize}
So the tactics keep the single term (\cref{sec:qed:assembly}), and
\code{MeBi Prove} is for the goals whose memory matters.

\subsection{Cost}\label{sec:qed:cost}

Measured on 11 November 2026 (the proof built by the tactic, then
\code{Qed}; memory at most, for the whole file; for \code{MeBi Prove},
the proof built with its pieces checked, then the lemma checked,
\cref{sec:qed:prove}):
\begin{center}
\begin{tabular}{lrrr}
\toprule
goal & built in & \code{Qed} & memory \\
\midrule
\file{LogTest.v}, each & 0.1\,s & -- & -- \\
Proc Test1--3, each & 0.1--0.4\,s & -- & -- \\
CADP, one process, each & 0.8\,s & 0.2\,s & 0.5\,GB \\
CADP, two processes (12,027 pairs) & 3.3\,s & 1.7\,s & 0.83\,GB \\
CADP, three processes (7.9M pairs) & 124\,s & 144\,s & 6.6\,GB \\
\quad by \code{MeBi Prove} (12 November) & 231\,s & 0.4\,s & 2.6\,GB \\
\bottomrule
\end{tabular}
\end{center}
Two processes took 104\,s, then 10\,s at \code{Qed}, 2.1\,GB, before the
measures of \cref{sec:qed:fit}; three processes did not fit in 7\,GB.
Now the three-process proof's main term holds 828,865 facts, and the
packs' memberships, about 280,000 facts in all, are checked as constants
while the proof is built (59\,s of the 124\,s). The whole file,
\file{examples/CADP/ThreeProcesses.v}, compiles in 293\,s, with 6.6\,GB
at most, the compiled file written (113\,MB); by \code{MeBi Prove},
in 279\,s, with 2.6\,GB at most (\cref{sec:qed:prove}), as the example
now does. Every example is therefore proved. What the cost grows with is the certificate itself:
the facts the kernel checks, about 6\,KB of its memory each, which only
the optimisations of \cref{sec:optimisations} can reduce.

\subsection{Trying it}\label{sec:qed:try}

Any example in \file{examples/} or \file{test/rocq/LogTest.v}: each goal
ends with \code{Qed}. To see the phases, set \code{MEBI_CERT_STATS=1}:
lines on standard error, each with the time since the start and the
process's memory, now and at most:
\begin{itemize}
\item \code{[qed] coverage of shape \dots: N projections}, then
  \code{coverage: N leaves built (\dots)}, then \code{coverage: declared}:
  a coverage's leaves, and its constant checked;
\item \code{chunk N written: F facts in all (lemma count, \dots)}: the
  translator's chunks (the counts add up), with the facts per lemma;
\item \code{[qed] kernel, per kind of constant: \dots}: the kernel's time
  per kind of piece (\code{case}, \code{as} for step lemmas, \code{lf} for
  groups of cases, \code{all} for coverages, \code{pack}, \code{snd} for
  soundness, \dots);
\item \code{[qed] before Qed}: the memory left for the kernel's check.
\end{itemize}
\code{MEBI_PLAN_DUMP=1} prints what the engine recorded (shapes, values,
relations and their tuples, summands and their projections);
\code{MEBI_QED_TEXT=<file>} appends every text the proof elaborates to
\code{<file>}; \code{MEBI_QED_LETS=1} makes the pieces \code{let}s of the
term instead of constants (to compare). Run large examples under a
memory cap (\file{CLAUDE.md}).
